
\documentclass{article}
\usepackage{colortbl}
\usepackage{makecell}
\usepackage{multirow}
\usepackage{supertabular}

\begin{document}

\newcounter{utterance}

\centering \large Interaction Transcript for game `cladder', experiment `full\_v1.5\_default', episode 9763 with llama{-}3.1{-}8b{-}instruct.
\vspace{24pt}

{ \footnotesize  \setcounter{utterance}{1}
\setlength{\tabcolsep}{0pt}
\begin{supertabular}{c@{$\;$}|p{.15\linewidth}@{}p{.15\linewidth}p{.15\linewidth}p{.15\linewidth}p{.15\linewidth}p{.15\linewidth}}
    \# & \multicolumn{2}{c}{Player} && \multicolumn{2}{c}{Game Master} \\
    \hline

    \theutterance \stepcounter{utterance}  
    & & & \multicolumn{4}{p{0.6\linewidth}}{
        \cellcolor[rgb]{0.9,0.9,0.9}{
            \makecell[{{p{\linewidth}}}]{
                \texttt{\tiny{[P1$\langle$GM]}}
                \texttt{You are an expert in causal inference. The following question is not a typical commonsense query, but rather a meticulously designed question created by a professor specializing in causal inference, intended to assess the students' mastery of the course content.} \\
\texttt{Imagine a self{-}contained, hypothetical world with only the following conditions, and without any unmentioned factors or causal relationships: Yupt has a direct effect on jyka and kwox. Jyka has a direct effect on kwox. We know that yupt causes jyka. yupt and jyka causes kwox. We observed an individual is yupt. Would an individual is kwox if not jyka instead of jyka?} \\
\texttt{Start your answer with 'Yes' or 'No', followed by additional reasoning or evidence to support your explanation.} \\
            }
        }
    }
     \\ \\

    \theutterance \stepcounter{utterance}  
    & \multicolumn{4}{p{0.6\linewidth}}{
        \cellcolor[rgb]{0.561,0.737,0.561}{
            \makecell[{{p{\linewidth}}}]{
                \texttt{\tiny{[P1$\rangle$GM]}}
                \texttt{No.} \\
\\ 
\texttt{To answer this question, we need to apply the principles of causal inference, specifically the concept of causal mediation and the backdoor criterion.} \\
\\ 
\texttt{Given the conditions:} \\
\\ 
\texttt{1. Yupt has a direct effect on jyka and kwox.} \\
\texttt{2. Jyka has a direct effect on kwox.} \\
\texttt{3. Yupt causes jyka.} \\
\texttt{4. Yupt and jyka cause kwox.} \\
\\ 
\texttt{We are asked to evaluate the effect of not having jyka (instead of jyka) on the outcome kwox, given that we observed an individual is yupt.} \\
\\ 
\texttt{To do this, we need to consider the causal pathway from yupt to kwox. There are two possible pathways:} \\
\\ 
\texttt{1. Direct effect of yupt on kwox (Yupt → Kwox)} \\
\texttt{2. Mediated effect of yupt on kwox through jyka (Yupt → Jyka → Kwox)} \\
\\ 
\texttt{Since we know that yupt causes jyka, and jyka causes kwox, the mediated effect (Yupt → Jyka → Kwox) is a part of the total effect of yupt on kwox.} \\
\\ 
\texttt{Now, if we were to remove jyka from the equation (i.e., not having jyka instead of jyka), we would be left with only the direct effect of yupt on kwox (Yupt → Kwox). However, we cannot conclude that this direct effect would be sufficient to cause kwox, as the direct effect of yupt on kwox is not specified.} \\
\\ 
\texttt{Moreover, the presence of the direct effect of yupt on kwox does not imply that the absence of jyka would not affect the outcome kwox. In fact, the backdoor criterion suggests that if we were to remove jyka, we would be introducing a new confounding variable (jyka) that could affect the outcome kwox.} \\
\\ 
\texttt{Therefore, based on the given conditions and the principles of causal inference, we cannot conclude that an individual would be kwox if not jyka instead of jyka, given that we observed an individual is yupt. The answer is No.} \\
            }
        }
    }
    & & \\ \\

    \theutterance \stepcounter{utterance}  
    & & & \multicolumn{2}{p{0.3\linewidth}}{
        \cellcolor[rgb]{0.9,0.9,0.9}{
            \makecell[{{p{\linewidth}}}]{
                \texttt{\tiny{[GM$|$GM]}}
                \texttt{No} \\
            }
        }
    }
    & & \\ \\

    \theutterance \stepcounter{utterance}  
    & & & \multicolumn{2}{p{0.3\linewidth}}{
        \cellcolor[rgb]{0.9,0.9,0.9}{
            \makecell[{{p{\linewidth}}}]{
                \texttt{\tiny{[GM$|$GM]}}
                \texttt{no} \\
            }
        }
    }
    & & \\ \\

    \theutterance \stepcounter{utterance}  
    & & & \multicolumn{2}{p{0.3\linewidth}}{
        \cellcolor[rgb]{0.9,0.9,0.9}{
            \makecell[{{p{\linewidth}}}]{
                \texttt{\tiny{[GM$|$GM]}}
                \texttt{game\_result = WIN} \\
            }
        }
    }
    & & \\ \\

\end{supertabular}
}

\end{document}
